\pagestyle{empty}
\tight
\begin{tblr}{width=\textwidth,colspec={|Q[c,m,1.9cm]|X[l,m]|},hlines,vlines,hspan=minimal,row{1}={bg=headblue},rowsep=1pt,colsep=3pt}
\SetCell{c,m}\hfil\logo\hfil & \textbf{POLITEKNIK NEGERI MALANG}\par\textbf{JURUSAN TEKNOLOGI INFORMASI}\par\textbf{PROGRAM STUDI: D-IV SISTEM INFORMASI BISNIS} \\
\SetCell[c=2]{c} \textbf{RENCANA PEMBELAJARAN SEMESTER (RPS)} & \\
\end{tblr}
\begin{longtblr}{width=\textwidth,colspec={|Q[l,m,1.9cm]|X[0.85,l,m]|X[1.45,l,m]|X[0.75,l,m]|X[1.15,l,m]|},hlines,vlines,hspan=minimal,rowsep=1pt,colsep=2pt,row{1,3}={bg=headgray,font=\bfseries}}
MATA KULIAH & KODE & BOBOT (SKS)/jam & SEMESTER & TGL. PENYUSUNAN \\
Pemrograman Web Lanjut & SIB245007 & 3 SKS/6 jam & 4 & 28 Juli 2026 \\
OTORISASI & Dosen Pengembang RPS & Koordinator RMK & \SetCell[c=2]{l} Ka PRODI &  \\
 & \begin{enumerate}\item Dian Hanifudin Subhi, S.Kom., M.Kom.\item Ade Ismail, S.Kom., M.TI.\end{enumerate} &  & \SetCell[c=2]{l}{\vspace{8mm}\par Hendra Pradibta, S.E., M.Sc.} &  \\
\SetCell[r=12]{l} Capaian Pembelajaran (CP) & \SetCell[c=4]{l,bg=headgray,font=\bfseries} Capaian Pembelajaran Lulusan Program Studi (CPL-Prodi) &  &  &  \\
 & CPL02 & \SetCell[c=3]{l} Mampu merancang, mengimplementasikan, dan menganalisis algoritma menggunakan berbagai paradigma pemrograman, serta mampu memecahkan masalah komputasional dengan logika yang sistematis dan tepat. &  &  \\
 & \SetCell[c=4]{l,bg=headgray,font=\bfseries} Capaian Pembelajaran mata kuliah (CPL-MK) &  &  &  \\
 & CPMK0202 & \SetCell[c=3]{l} Mahasiswa mampu mengimplementasikan algoritma dalam program komputer menggunakan berbagai paradigma pemrograman, seperti prosedural, berorientasi objek, dan berbasis web/mobile. &  &  \\
 & \SetCell[c=4]{l,bg=headgray,font=\bfseries} Kemampuan akhir tiap tahapan belajar (Sub-CPMK) &  &  &  \\
 & SCPMK0202-04001 & \SetCell[c=3]{l} Mahasiswa mampu memahami konsep dasar web framework serta menerapkan routing, controller, dan pengelolaan basis data dalam pengembangan aplikasi web. &  &  \\
 & SCPMK0202-04002 & \SetCell[c=3]{l} Mahasiswa mampu menerapkan model, server-side templating, dan operasi CRUD dalam pengembangan aplikasi web berbasis framework. &  &  \\
 & SCPMK0202-04003 & \SetCell[c=3]{l} Mahasiswa mampu menerapkan mekanisme autentikasi, otorisasi berbasis peran, serta pengembangan dan konsumsi API pada aplikasi web. &  &  \\
 & SCPMK0202-04004 & \SetCell[c=3]{l} Mahasiswa mampu mengembangkan aplikasi web berbasis framework secara terintegrasi sesuai dengan kebutuhan pengguna. &  &  \\
\SetCell[r=2]{l} Penilaian & \SetCell[c=4]{l} *Nilai CPMK didapat dari agregasi nilai SUB-CPMK yang dibebankan &  &  &  \\
 & \SetCell[c=4]{l}{\tiny\begin{tblr}{width=\linewidth,colspec={|Q[c,m,0.11\linewidth]|Q[c,m,0.18\linewidth]|X[l,m]|X[l,m]|X[l,m]|X[l,m]|Q[c,m,0.08\linewidth]|},hlines,vlines,hspan=minimal,rowsep=1pt,colsep=0.7pt,row{1,2}={bg=headgray,font=\bfseries}}
Kode CPMK & Kode SCPMK & \SetCell[c=4]{c} Bobot per Bentuk Penilaian & & & & Total Bobot \\
 & & Tugas & UTS & PBL & UAS & \\
CPMK0202 & SCPMK0202-04001 & 5\%: setup framework, routing, controller, pengelolaan basis data & 2\%: evaluasi progress proyek & & & 7\% \\
CPMK0202 & SCPMK0202-04002 & 5\%: templating, ORM, CRUD, validasi input & 2\%: evaluasi progress proyek & & & 7\% \\
CPMK0202 & SCPMK0202-04003 & 3\%: autentikasi, otorisasi, desain \& konsumsi API & 1\%: evaluasi progress proyek & & & 4\% \\
CPMK0202 & SCPMK0202-04004 & 2\%: pengolahan data \& import/export & & 55\%: proyek web terintegrasi (proposal minggu 4 \& pengembangan berkelanjutan) & 25\%: proyek final & 82\% \\
\SetCell[c=6]{r}\textbf{Total Per Penilaian} & & & & & & \textbf{100\%} \\
\end{tblr}} &  &  &  \\
Diskripsi Singkat Mata Kuliah & \SetCell[c=4]{l} Mata kuliah Pemrograman Web Lanjut membahas pengembangan aplikasi web modern menggunakan pendekatan framework berbasis MVC, mencakup routing, controller, pengelolaan basis data, ORM, server-side templating, validasi input, autentikasi, otorisasi, pengolahan data, desain API, serta pengembangan proyek web terintegrasi. Pembelajaran berjalan secara inkremental melalui studi kasus aplikasi Point of Sale (POS)/kasir UMKM yang dibangun bertahap tiap minggu (minggu 1--10) sekaligus menjadi rujukan arsitektur bagi proyek PBL kelompok yang berjalan sepanjang semester: proposal proyek diajukan pada minggu 4, pengembangan berlanjut hingga minggu 16, dan proyek final dievaluasi pada UAS. Konsep bersifat umum dan dapat diterapkan pada berbagai framework web (mis. Laravel, Django, Express, Spring). &  &  &  \\
Bahan Kajian / Materi Pembelajaran & \SetCell[c=4]{l}\begin{enumerate}\item Arsitektur web modern, pengantar framework MVC, routing, controller, dan pengelolaan basis data\item Server-side templating, ORM, validasi input, dan operasi CRUD\item Autentikasi, otorisasi berbasis peran (RBAC), pengolahan data, desain API, dan proyek web terintegrasi\end{enumerate} &  &  &  \\
\SetCell[r=2]{l} Pustaka & \SetCell[c=4]{l} \textbf{Utama:}\begin{enumerate}\item Welling, L. and Thomson, L. 2017. PHP and MySQL Web Development, 5th ed. Addison-Wesley.\item Lockhart, J. 2015. Modern PHP: New Features and Good Practices. O'Reilly Media.\end{enumerate} &  &  &  \\
 & \SetCell[c=4]{l} \textbf{Pendukung:} Materi LMS, dokumentasi resmi framework web yang digunakan (mis. laravel.com/docs), dan jobsheet praktikum. &  &  &  \\
Media Pembelajaran & \SetCell[c=2]{l} \textbf{Software:} Framework web, RDBMS, dependency manager, version control, code editor (contoh: PHP/Laravel, MySQL, Composer, Git, VS Code). &  & \SetCell[c=2]{l} \textbf{Hardware:} PC/laptop, proyektor. &  \\
Nama Dosen Pengampu & \SetCell[c=4]{l}\begin{enumerate}\item Dian Hanifudin Subhi, S.Kom., M.Kom.\item Ade Ismail, S.Kom., M.TI.\end{enumerate} &  &  &  \\
Matakuliah Syarat & \SetCell[c=4]{l} Desain dan Pemrograman Web &  &  &  \\
\end{longtblr}
\begin{longtblr}{width=\textwidth,colspec={|Q[c,m,0.06\textwidth]|X[1.35,l,m]|X[1.08,l,m]|X[1.16,l,m]|X[0.7,l,m]|X[1.24,l,m]|X[1.06,l,m]|X[0.98,l,m]|Q[c,m,0.05\textwidth]|},hlines,vlines,hspan=minimal,rowhead=2,row{1,2}={bg=headgreen,font=\bfseries},rowsep=1pt,colsep=1.5pt}
Mgg.\par Ke & Kemampuan Akhir Yang Direncanakan (Sub-CP-MK) & Bahan kajian (Materi Pembelajaran) & Bentuk dan Metode Pembelajaran & Estimasi Waktu & Pengalaman Belajar Mahasiswa & Kriteria \& Bentuk Penilaian & Indikator Penilaian & Bobot Penilaian (\%) \\
(1) & (2) & (3) & (4) & (5) & (6) & (7) & (8) & (9) \\
1 & SCPMK0202-04001 & Arsitektur web modern (monolith, microservices, serverless); pengantar web framework; perbandingan framework; lingkungan pengembangan; version control dasar. Studi kasus: pengenalan \& setup proyek aplikasi POS/kasir UMKM, inisialisasi repository. & Modalitas: Blended Learning. Bentuk: Kuliah \& Praktikum. Metode: small group discussion, hands-on practice. & 1 x 2 x 50' tatap muka; 1 x 2 x 50' tugas/praktik mandiri. & Mahasiswa mendiskusikan evolusi arsitektur web, menyiapkan lingkungan pengembangan, dan menginisialisasi proyek studi kasus aplikasi POS/kasir UMKM (increment 1); kelompok PBL mulai membentuk tim dan menjajaki domain bisnis. & Bentuk penilaian: Tugas individu, setup lingkungan \& inisialisasi proyek POS; kontribusi awal proyek PBL (pembentukan tim). Mengacu rubrik RTM. & Ketepatan konsep; kualitas artefak/praktik; validasi dan dokumentasi. & 4 \\
2 & SCPMK0202-04001 & HTTP protocol (request/response, methods, status codes, headers); arsitektur MVC; routing; organisasi controller. Studi kasus POS: route \& controller halaman produk dan transaksi kasir. & Modalitas: Blended Learning. Bentuk: Kuliah \& Praktikum. Metode: problem-based learning, hands-on practice. & 1 x 2 x 50' tatap muka; 1 x 2 x 50' tugas/praktik mandiri. & Mahasiswa mengimplementasikan routing dan controller produk \& transaksi kasir pada aplikasi studi kasus POS (increment 2), serta menunjukkan evidence hasil belajar; kelompok PBL melanjutkan eksplorasi domain bisnis. & Bentuk penilaian: Tugas individu, routing \& controller POS; kontribusi awal proyek PBL (eksplorasi domain). Mengacu rubrik RTM. & Ketepatan konsep; kualitas artefak/praktik; validasi dan dokumentasi. & 4 \\
3 & SCPMK0202-04002 & Server-side templating \& komponen tampilan (view components); asset bundling \& build tooling frontend; utility-first CSS; perbandingan MPA vs SPA. Studi kasus POS: layout daftar produk \& form kasir. & Modalitas: Blended Learning. Bentuk: Kuliah \& Praktikum. Metode: problem-based learning, hands-on practice. & 1 x 2 x 50' tatap muka; 1 x 2 x 50' tugas/praktik mandiri. & Mahasiswa membangun antarmuka daftar produk dan form kasir aplikasi POS (increment 3) menggunakan server-side templating dan komponen tampilan, serta menunjukkan evidence hasil belajar; kelompok PBL menuntaskan pemilihan domain bisnis sebagai persiapan proposal minggu depan. & Bentuk penilaian: Tugas individu, templating \& komponen tampilan POS; kontribusi awal proyek PBL (finalisasi domain). Mengacu rubrik RTM. & Ketepatan konsep; kualitas artefak/praktik; validasi dan dokumentasi. & 4 \\
4 & SCPMK0202-04001, SCPMK0202-04004 & Desain basis data; strategi indexing; migrasi skema (migration, schema builder, seeding, factory); pengantar ORM; perbandingan ORM lintas framework. Studi kasus POS: skema produk, kategori, transaksi, detail transaksi. Proyek PBL: pengajuan proposal proyek (domain bisnis pilihan tim, konsep tingkat tinggi, daftar fitur awal). & Modalitas: Blended Learning. Bentuk: Kuliah \& Praktikum \& Proyek/PBL. Metode: problem-based learning, hands-on practice, project-based learning. & 1 x 2 x 50' tatap muka; 1 x 2 x 50' tugas/praktik mandiri. & Mahasiswa merancang skema basis data POS (produk, kategori, transaksi, detail transaksi) dan mengimplementasikan migrasi serta seeding (increment 4); kelompok PBL mengajukan proposal proyek yang memuat domain bisnis pilihan, konsep aplikasi, dan daftar fitur awal, dengan aplikasi studi kasus POS sebagai referensi arsitektur. & Bentuk penilaian: Tugas individu, desain basis data \& migrasi POS; Tugas kelompok, proposal proyek PBL. Mengacu rubrik RTM. & Ketepatan konsep; kualitas artefak/praktik; validasi dan dokumentasi; kesesuaian dan kelengkapan proposal. & 16 \\
5 & SCPMK0202-04002 & ORM \& pola akses data; relasi antar model (one-to-one, one-to-many, many-to-many); accessor/mutator; eager loading. Studi kasus POS: relasi transaksi--detail transaksi, produk--kategori, kasir--transaksi. & Modalitas: Blended Learning. Bentuk: Kuliah \& Praktikum. Metode: problem-based learning, hands-on practice. & 1 x 2 x 50' tatap muka; 1 x 2 x 50' tugas/praktik mandiri. & Mahasiswa mengimplementasikan relasi basis data POS (transaksi--detail, produk--kategori, kasir--transaksi) melalui ORM (increment 5), lalu menunjukkan evidence hasil belajar; kelompok PBL melanjutkan pengembangan proyek berdasarkan proposal yang telah disetujui. & Bentuk penilaian: Tugas individu, ORM \& relasi data POS; kontribusi progress proyek PBL. Mengacu rubrik RTM. & Ketepatan konsep; kualitas artefak/praktik; validasi dan dokumentasi. & 4 \\
6 & SCPMK0202-04002 & Validasi \& sanitasi input; validasi berbasis request/form; flash message; penanganan error \& UX. Studi kasus POS: validasi form produk (stok, harga) \& input kasir (qty). & Modalitas: Blended Learning. Bentuk: Kuliah \& Praktikum. Metode: problem-based learning, hands-on practice. & 1 x 2 x 50' tatap muka; 1 x 2 x 50' tugas/praktik mandiri. & Mahasiswa mengimplementasikan validasi \& sanitasi input pada form produk dan transaksi kasir aplikasi POS (increment 6), serta menunjukkan evidence hasil belajar; kelompok PBL melanjutkan pengembangan proyek. & Bentuk penilaian: Tugas individu, validasi \& keamanan input POS; kontribusi progress proyek PBL. Mengacu rubrik RTM. & Ketepatan konsep; kualitas artefak/praktik; validasi dan dokumentasi. & 3 \\
7 & SCPMK0202-04003 & Autentikasi (authentication scaffolding); otorisasi berbasis peran (RBAC); manajemen hak akses (role \& permission management); perbandingan pendekatan autentikasi lintas framework. Studi kasus POS: login \& peran admin/kasir. & Modalitas: Blended Learning. Bentuk: Kuliah \& Praktikum. Metode: problem-based learning, hands-on practice. & 1 x 2 x 50' tatap muka; 1 x 2 x 50' tugas/praktik mandiri. & Mahasiswa mengimplementasikan autentikasi dan otorisasi peran admin/kasir pada aplikasi POS (increment 7), serta menunjukkan evidence hasil belajar; kelompok PBL melanjutkan pengembangan proyek. & Bentuk penilaian: Tugas individu, autentikasi \& otorisasi POS; kontribusi progress proyek PBL. Mengacu rubrik RTM. & Ketepatan konsep; kualitas artefak/praktik; validasi dan dokumentasi. & 4.5 \\
8 & SCPMK0202-04001, SCPMK0202-04002, SCPMK0202-04003 & UTS: Evaluasi milestone aplikasi studi kasus POS/kasir UMKM (increment 1--7), CRUD lengkap; REST principles; arsitektur MVC; desain basis data; validasi \& autentikasi; defense lisan (viva). & Modalitas: Tatap Muka. Bentuk: Ujian Praktik. Metode: hands-on practice, viva, presentasi. & 1 x 2 x 50' tatap muka; 1 x 2 x 50' praktik mandiri. & Mahasiswa mendemonstrasikan milestone aplikasi studi kasus POS yang mencakup CRUD produk/transaksi, autentikasi, dan arsitektur MVC, kemudian mempertahankan keputusan teknis secara lisan. & Bentuk penilaian: UTS Praktik, milestone aplikasi studi kasus \& defense. Mengacu rubrik RTM. & Kelengkapan fitur; kualitas arsitektur; kemampuan defense teknis. & 5 \\
9 & SCPMK0202-04004 & Pengolahan data (ETL, batch processing, antrian); manajemen memori; import/export data (library spreadsheet \& PDF generation); penanganan dataset besar. Studi kasus POS: laporan penjualan (export) \& import data produk. & Modalitas: Blended Learning. Bentuk: Kuliah \& Praktikum. Metode: problem-based learning, hands-on practice. & 1 x 2 x 50' tatap muka; 1 x 2 x 50' tugas/praktik mandiri. & Mahasiswa mengimplementasikan fitur laporan penjualan (export) dan import data produk pada aplikasi POS (increment 8), serta menunjukkan evidence hasil belajar; kelompok PBL melanjutkan pengembangan proyek. & Bentuk penilaian: Tugas individu, pengolahan \& import/export data POS; kontribusi progress proyek PBL. Mengacu rubrik RTM. & Ketepatan konsep; kualitas artefak/praktik; validasi dan dokumentasi. & 5 \\
10 & SCPMK0202-04003 & Arsitektur \& desain API; spesifikasi OpenAPI/Swagger; transformasi respons API (API resource); token-based API authentication; konsumsi API dari sisi klien (pemanggilan endpoint terproteksi, penanganan token dan galat respons). Studi kasus POS: API produk \& transaksi untuk kasir mobile, serta contoh klien yang mengonsumsinya. & Modalitas: Blended Learning. Bentuk: Kuliah \& Praktikum. Metode: problem-based learning, hands-on practice. & 1 x 2 x 50' tatap muka; 1 x 2 x 50' tugas/praktik mandiri. & Mahasiswa merancang dan mengimplementasikan endpoint API produk \& transaksi untuk kasir mobile pada aplikasi POS (increment 9, penutup studi kasus) berikut dokumentasinya, membangun contoh klien yang mengonsumsi API tersebut (login mendapatkan token, memanggil endpoint terproteksi, menangani respons dan galat), serta menunjukkan evidence hasil belajar; kelompok PBL melanjutkan pengembangan proyek. & Bentuk penilaian: Tugas individu, desain, implementasi, \& konsumsi API POS; kontribusi progress proyek PBL. Mengacu rubrik RTM. & Ketepatan konsep; kualitas artefak/praktik; validasi dan dokumentasi. & 4 \\
11 & SCPMK0202-04004 & PBL proyek web, perencanaan teknis lanjutan; arsitektur aplikasi terperinci; pembagian peran tim. Tim mendetailkan proposal yang sudah disetujui pada minggu 4 (domain bisnis, mis. inventori, koperasi, reservasi, laundry) menjadi rancangan arsitektur dan rencana eksekusi konkret, dengan aplikasi studi kasus POS sebagai referensi arsitektur dan inspirasi fitur. & Modalitas: Blended Learning. Bentuk: Proyek/PBL. Metode: project-based learning, mentoring. & 1 x 2 x 50' tatap muka; 1 x 2 x 50' praktik mandiri. & Mahasiswa mendetailkan proposal PBL menjadi arsitektur dan rencana eksekusi konkret dengan mengacu pada pola aplikasi studi kasus POS, membagi peran tim, kemudian mendokumentasikan rencana pengembangan. & Bentuk penilaian: Tugas kelompok, perencanaan teknis proyek. Mengacu rubrik RTM. & Kesesuaian rencana; kelengkapan dokumentasi; peran tim. & 3 \\
12 & SCPMK0202-04004 & PBL proyek web, pengembangan fitur utama; code review sesama tim; mengadaptasi pola studi kasus POS ke domain tim. & Modalitas: Blended Learning. Bentuk: Proyek/PBL. Metode: project-based learning, mentoring, code review. & 1 x 2 x 50' tatap muka; 1 x 2 x 50' praktik mandiri. & Mahasiswa membangun fitur utama proyek web dengan mengadaptasi pola studi kasus POS ke domain tim, dan melakukan code review bersama rekan tim. & Bentuk penilaian: Tugas kelompok, fitur utama \& code review. Mengacu rubrik RTM. & Kualitas fitur; konsistensi kode; kolaborasi tim. & 3 \\
13 & SCPMK0202-04004 & PBL proyek web, integrasi modul; pengujian fungsional; perbaikan iteratif. & Modalitas: Blended Learning. Bentuk: Proyek/PBL. Metode: project-based learning, mentoring. & 1 x 2 x 50' tatap muka; 1 x 2 x 50' praktik mandiri. & Mahasiswa mengintegrasikan modul-modul proyek, melakukan pengujian fungsional, dan memperbaiki temuan. & Bentuk penilaian: Tugas kelompok, integrasi \& pengujian. Mengacu rubrik RTM. & Kelengkapan integrasi; cakupan pengujian; perbaikan temuan. & 3 \\
14 & SCPMK0202-04004 & PBL proyek web, optimasi performa; persiapan deployment; dokumentasi teknis. & Modalitas: Blended Learning. Bentuk: Proyek/PBL. Metode: project-based learning, mentoring. & 1 x 2 x 50' tatap muka; 1 x 2 x 50' praktik mandiri. & Mahasiswa mengoptimalkan performa aplikasi dan menyiapkan artefak deployment serta dokumentasi teknis. & Bentuk penilaian: Tugas kelompok, optimasi \& deployment preparation. Mengacu rubrik RTM. & Performa aplikasi; kelengkapan dokumentasi; kesiapan deployment. & 4 \\
15 & SCPMK0202-04004 & PBL proyek web, finalisasi aplikasi web utuh; dokumentasi arsitektur; video demo; presentasi. & Modalitas: Blended Learning. Bentuk: Proyek/PBL. Metode: project-based learning, mentoring, presentasi. & 1 x 2 x 50' tatap muka; 1 x 2 x 50' praktik mandiri. & Mahasiswa menyelesaikan dan mempresentasikan proyek web terintegrasi beserta dokumentasinya kepada dosen/reviewer. & Bentuk penilaian: Tugas kelompok, finalisasi \& presentasi proyek. Mengacu rubrik RTM. & Kelengkapan fitur; kualitas dokumentasi; presentasi proyek. & 4 \\
16 & SCPMK0202-04004 & PBL proyek web, persiapan ujian akhir; sesi rehearsal/mock defense. Tiap kelompok mempresentasikan kondisi proyek terkini di hadapan mentor/rekan sejawat, menerima umpan balik, dan menuntaskan perbaikan sebelum UAS sesungguhnya. & Modalitas: Blended Learning. Bentuk: Proyek/PBL. Metode: project-based learning, mentoring, mock presentation. & 1 x 2 x 50' tatap muka; 1 x 2 x 50' praktik mandiri. & Mahasiswa mempresentasikan kondisi proyek terkini secara simulasi, menerima umpan balik dari mentor/rekan sejawat, dan menindaklanjuti perbaikan sebelum presentasi UAS. & Bentuk penilaian: Tugas kelompok, rehearsal/mock defense. Mengacu rubrik RTM. & Kesiapan presentasi; responsivitas terhadap umpan balik; kematangan proyek. & 4.5 \\
17 & SCPMK0202-04004 & UAS: Evaluasi proyek web final, presentasi produk; defense lisan (viva); evaluasi ketercapaian kebutuhan pengguna. & Modalitas: Tatap Muka. Bentuk: Ujian Praktik. Metode: viva, presentasi, evaluasi. & 1 x 2 x 50' tatap muka; 1 x 2 x 50' praktik mandiri. & Mahasiswa mendemonstrasikan proyek web final, mempertahankan keputusan teknis secara lisan, dan mengevaluasi ketercapaian kebutuhan pengguna. & Bentuk penilaian: UAS Praktik, proyek final \& defense. Mengacu rubrik RTM. & Kesesuaian kebutuhan pengguna; kualitas teknis; kemampuan defense. & 25 \\
\end{longtblr}
